% GENERATED FILE. Edit canonical lessons, then run npm run generate:books.
% canonical-source-hash: fnv1a64:24b4cbe9db322204
% canonical-lessons: SA-C401-bhiksuh, SA-C401-viharah, SA-C401-bhiksa, SA-C401-upasakah, SA-C401-caityam, SA-R401-review-394-401

\chapter{Monks and Monasteries}
\label{ch:sa-monks-and-monasteries}
\hlchaptermodality{401}

\begin{chapteropening}
\textbf{By the end of this chapter:} \emph{I can talk about monks, alms and the shrines they keep.}

Complete the last lesson of chapter 401: I can talk about monks, alms and the shrines they keep.
\end{chapteropening}

\section[\texorpdfstring{\sk{भिक्षुः} (bhikṣuḥ)}{भिक्षुः (bhikṣuḥ)}]{\sk{भिक्षुः} (bhikṣuḥ) — a mendicant monk}

\label{lesson:SA-C401-bhiksuh}



\begin{quote}
Before the new one: say the Sanskrit for a chariot procession, then the Sanskrit for fetches.
\end{quote}

\subsection*{You'll want to know: \sk{भिक्षुः}}
\textbf{\sk{भिक्षुः}} — \emph{bhikṣuḥ} — \textquotedblleft{}a mendicant monk\textquotedblright{}.

\begin{quote}
\sk{भिक्षुः} \sk{प्रातः} \sk{ग्रामं} \sk{गच्छति।}
\end{quote}

\begin{quote}
\emph{bhikṣuḥ prātaḥ grāmaṃ gacchati.}
\end{quote}

\begin{quote}
The monk goes to the village in the morning.
\end{quote}

A masculine noun.

\begin{cousinweb}
\textbf{\sk{भिक्षुः}} (\emph{bhikṣuḥ}) comes from the root √bhikṣ, \textquotedblleft{}to beg, to ask for alms\textquotedblright{}; MW gives \emph{bhikṣu} as a beggar, a religious mendicant and a Buddhist monk (MW).
\end{cousinweb}

\begin{grammarlens}[title={What you've built}]
One more word for getting around.
\end{grammarlens}

\subsection*{Your turn}
\begin{itemize}
\raggedright
  \item \emph{Say it:} \emph{bhikṣuḥ}
  \item \emph{Say it:} \emph{bhikṣuḥ}, once more
  \item \emph{Say it:} \emph{bhikṣuḥ prātaḥ grāmaṃ gacchati.} — The monk goes to the village in the morning.
  \item \emph{Recall:} say the Sanskrit for a chariot procession, then the Sanskrit for fetches, then say \emph{bhikṣuḥ} again
\end{itemize}

\begin{tcolorbox}[breakable,colback=teal!4,colframe=teal!35!black,title={Before you move on}]
What does \sk{भिक्षुः} mean? (\textquotedblleft{}A mendicant monk\textquotedblright{}.) Say it once more.
\end{tcolorbox}
\section[\texorpdfstring{\sk{विहारः} (vihāraḥ)}{विहारः (vihāraḥ)}]{\sk{विहारः} (vihāraḥ) — a Buddhist monastery}

\label{lesson:SA-C401-viharah}



\begin{quote}
Before the new one: say the Sanskrit for fetches, then the Sanskrit for a mendicant monk.
\end{quote}

\subsection*{You'll want to know: \sk{विहारः}}
\textbf{\sk{विहारः}} — \emph{vihāraḥ} — \textquotedblleft{}a Buddhist monastery\textquotedblright{}.

\begin{quote}
\sk{भिक्षुः} \sk{विहारे} \sk{वसति।}
\end{quote}

\begin{quote}
\emph{bhikṣuḥ vihāre vasati.}
\end{quote}

\begin{quote}
The monk lives in the monastery.
\end{quote}

A masculine noun.

\begin{cousinweb}
\textbf{\sk{विहारः}} (\emph{vihāraḥ}) is built from \emph{vi}, \textquotedblleft{}apart\textquotedblright{}, and the root √hṛ, \textquotedblleft{}to carry\textquotedblright{}: MW gives \emph{vihāra} as walking for pleasure and, among Buddhists and Jainas, a monastery (MW).
\end{cousinweb}

\begin{grammarlens}[title={What you've built}]
One more word for getting around.
\end{grammarlens}

\subsection*{Your turn}
\begin{itemize}
\raggedright
  \item \emph{Say it:} \emph{vihāraḥ}
  \item \emph{Say it:} \emph{vihāraḥ}, once more
  \item \emph{Say it:} \emph{bhikṣuḥ vihāre vasati.} — The monk lives in the monastery.
  \item \emph{Recall:} say the Sanskrit for fetches, then the Sanskrit for a mendicant monk, then say \emph{vihāraḥ} again
\end{itemize}

\begin{tcolorbox}[breakable,colback=teal!4,colframe=teal!35!black,title={Before you move on}]
What does \sk{विहारः} mean? (\textquotedblleft{}A Buddhist monastery\textquotedblright{}.) Say it once more.
\end{tcolorbox}
\section[\texorpdfstring{\sk{भिक्षा} (bhikṣā)}{भिक्षा (bhikṣā)}]{\sk{भिक्षा} (bhikṣā) — alms}

\label{lesson:SA-C401-bhiksa}



\begin{quote}
Before the new one: say the Sanskrit for a mendicant monk, then the Sanskrit for a Buddhist monastery.
\end{quote}

\subsection*{You'll want to know: \sk{भिक्षा}}
\textbf{\sk{भिक्षा}} — \emph{bhikṣā} — \textquotedblleft{}alms\textquotedblright{}.

\begin{quote}
\sk{भिक्षुः} \sk{गृहे} \sk{गृहे} \sk{भिक्षां} \sk{याचते।}
\end{quote}

\begin{quote}
\emph{bhikṣuḥ gṛhe gṛhe bhikṣāṃ yācate.}
\end{quote}

\begin{quote}
The monk asks for alms from house to house.
\end{quote}

A feminine noun.

\begin{cousinweb}
\textbf{\sk{भिक्षा}} (\emph{bhikṣā}) comes from the root √bhikṣ, \textquotedblleft{}to beg\textquotedblright{}; MW gives \emph{bhikṣā} as begging and as the alms or food a beggar receives (MW).
\end{cousinweb}

\begin{grammarlens}[title={What you've built}]
One more word for getting around.
\end{grammarlens}

\subsection*{Your turn}
\begin{itemize}
\raggedright
  \item \emph{Say it:} \emph{bhikṣā}
  \item \emph{Say it:} \emph{bhikṣā}, once more
  \item \emph{Say it:} \emph{bhikṣuḥ gṛhe gṛhe bhikṣāṃ yācate.} — The monk asks for alms from house to house.
  \item \emph{Recall:} say the Sanskrit for a mendicant monk, then the Sanskrit for a Buddhist monastery, then say \emph{bhikṣā} again
\end{itemize}

\begin{tcolorbox}[breakable,colback=teal!4,colframe=teal!35!black,title={Before you move on}]
What does \sk{भिक्षा} mean? (\textquotedblleft{}Alms\textquotedblright{}.) Say it once more.
\end{tcolorbox}
\section[\texorpdfstring{\sk{उपासकः} (upāsakaḥ)}{उपासकः (upāsakaḥ)}]{\sk{उपासकः} (upāsakaḥ) — a lay follower}

\label{lesson:SA-C401-upasakah}



\begin{quote}
Before the new one: say the Sanskrit for a Buddhist monastery, then the Sanskrit for alms.
\end{quote}

\subsection*{You'll want to know: \sk{उपासकः}}
\textbf{\sk{उपासकः}} — \emph{upāsakaḥ} — \textquotedblleft{}a lay follower\textquotedblright{}.

\begin{quote}
\sk{उपासकः} \sk{भिक्षवे} \sk{भोजनं} \sk{ददाति।}
\end{quote}

\begin{quote}
\emph{upāsakaḥ bhikṣave bhojanaṃ dadāti.}
\end{quote}

\begin{quote}
The lay follower gives food to the monk.
\end{quote}

A masculine noun.

\begin{cousinweb}
\textbf{\sk{उपासकः}} (\emph{upāsakaḥ}) is built from \emph{upa}, \textquotedblleft{}near\textquotedblright{}, and the root √ās, \textquotedblleft{}to sit\textquotedblright{}, which also gives \emph{āste}, \textquotedblleft{}sits\textquotedblright{}: one who sits near and serves; MW gives \emph{upāsaka} as a servant, a worshipper and a follower (MW).
\end{cousinweb}

\begin{grammarlens}[title={What you've built}]
One more word for getting around.
\end{grammarlens}

\subsection*{Your turn}
\begin{itemize}
\raggedright
  \item \emph{Say it:} \emph{upāsakaḥ}
  \item \emph{Say it:} \emph{upāsakaḥ}, once more
  \item \emph{Say it:} \emph{upāsakaḥ bhikṣave bhojanaṃ dadāti.} — The lay follower gives food to the monk.
  \item \emph{Recall:} say the Sanskrit for a Buddhist monastery, then the Sanskrit for alms, then say \emph{upāsakaḥ} again
\end{itemize}

\begin{tcolorbox}[breakable,colback=teal!4,colframe=teal!35!black,title={Before you move on}]
What does \sk{उपासकः} mean? (\textquotedblleft{}A lay follower\textquotedblright{}.) Say it once more.
\end{tcolorbox}
\section[\texorpdfstring{\sk{चैत्यम्} (caityam)}{चैत्यम् (caityam)}]{\sk{चैत्यम्} (caityam) — a memorial shrine; a sacred mound}

\label{lesson:SA-C401-caityam}



\begin{quote}
Before the new one: say the Sanskrit for alms, then the Sanskrit for a lay follower.
\end{quote}

\subsection*{You'll want to know: \sk{चैत्यम्}}
\textbf{\sk{चैत्यम्}} — \emph{caityam} — \textquotedblleft{}a memorial shrine; a sacred mound\textquotedblright{}.

\begin{quote}
\sk{विहारस्य} \sk{समीपे} \sk{प्राचीनं} \sk{चैत्यम्} \sk{अस्ति।}
\end{quote}

\begin{quote}
\emph{vihārasya samīpe prācīnaṃ caityam asti.}
\end{quote}

\begin{quote}
Near the monastery there is an ancient shrine.
\end{quote}

A neuter noun.

\begin{cousinweb}
\textbf{\sk{चैत्यम्}} (\emph{caityam}) is built on \emph{citā}, \textquotedblleft{}a funeral pile\textquotedblright{}, or \emph{citi}, \textquotedblleft{}a heap\textquotedblright{}, as MW gives it: MW gives \emph{caitya} as a funeral monument or stūpa, a sacred tree on a mound, and a temple or place of worship (MW).
\end{cousinweb}

\begin{grammarlens}[title={What you've built}]
One more word for getting around.
\end{grammarlens}

\subsection*{Your turn}
\begin{itemize}
\raggedright
  \item \emph{Say it:} \emph{caityam}
  \item \emph{Say it:} \emph{caityam}, once more
  \item \emph{Say it:} \emph{vihārasya samīpe prācīnaṃ caityam asti.} — Near the monastery there is an ancient shrine.
  \item \emph{Recall:} say the Sanskrit for alms, then the Sanskrit for a lay follower, then say \emph{caityam} again
\end{itemize}

\begin{tcolorbox}[breakable,colback=teal!4,colframe=teal!35!black,title={Before you move on}]
What does \sk{चैत्यम्} mean? (\textquotedblleft{}A memorial shrine; a sacred mound\textquotedblright{}.) Say it once more.
\end{tcolorbox}
\section[(dialogue)]{Looking back: from Setting Out on Pilgrimage to Monks and Monasteries}

\label{lesson:SA-R401-review-394-401}



\begin{quote}
Say the Sanskrit for a pilgrimage and starts on a journey.
\end{quote}

\subsection*{Your turn}
Say these aloud:

\begin{itemize}
\raggedright
  \item a pilgrimage, starts on a journey, food for the journey, travels, an ascetic's water-pot
  \item a caravan, a caravan leader, goes along with, takes leave, a wandering ascetic
  \item rings out, a sight of the deity, an image of a deity, an inner sanctum, a devotee
  \item kindles, a lamp's flame, camphor, a hymn of praise, sings the praises of
  \item a bathing ghat, plunges into, a confluence, purifying, sprinkles over
  \item a deep pool, gives a bath to, divine, an arched gateway, a pennant
  \item walks around, prayer beads, a rudraksha bead, a chariot procession, fetches
  \item a mendicant monk, a Buddhist monastery, alms, a lay follower, a memorial shrine
\end{itemize}

\begin{tcolorbox}[breakable,colback=teal!4,colframe=teal!35!black,title={Before you move on}]
A pilgrimage? (\textbf{\sk{तीर्थयात्रा}}, \emph{tīrthayātrā}.) Starts on a journey? (\textbf{\sk{प्रतिष्ठते}}, \emph{pratiṣṭhate}.) Food for the journey? (\textbf{\sk{पाथेयम्}}, \emph{pātheyam}.) Travels? (\textbf{\sk{याति}}, \emph{yāti}.) An ascetic's water-pot? (\textbf{\sk{कमण्डलुः}}, \emph{kamaṇḍaluḥ}.) A caravan? (\textbf{\sk{सार्थः}}, \emph{sārthaḥ}.) A caravan leader? (\textbf{\sk{सार्थवाहः}}, \emph{sārthavāhaḥ}.) Goes along with? (\textbf{\sk{अनुयाति}}, \emph{anuyāti}.) Takes leave? (\textbf{\sk{आपृच्छते}}, \emph{āpṛcchate}.) A wandering ascetic? (\textbf{\sk{परिव्राजकः}}, \emph{parivrājakaḥ}.) Rings out? (\textbf{\sk{ध्वनति}}, \emph{dhvanati}.) A sight of the deity? (\textbf{\sk{दर्शनम्}}, \emph{darśanam}.) An image of a deity? (\textbf{\sk{प्रतिमा}}, \emph{pratimā}.) An inner sanctum? (\textbf{\sk{गर्भगृहम्}}, \emph{garbhagṛham}.) A devotee? (\textbf{\sk{भक्तः}}, \emph{bhaktaḥ}.) Kindles? (\textbf{\sk{ज्वालयति}}, \emph{jvālayati}.) A lamp's flame? (\textbf{\sk{दीपशिखा}}, \emph{dīpaśikhā}.) Camphor? (\textbf{\sk{कर्पूरः}}, \emph{karpūraḥ}.) A hymn of praise? (\textbf{\sk{स्तोत्रम्}}, \emph{stotram}.) Sings the praises of? (\textbf{\sk{स्तौति}}, \emph{stauti}.) A bathing ghat? (\textbf{\sk{घट्टः}}, \emph{ghaṭṭaḥ}.) Plunges into? (\textbf{\sk{अवगाहते}}, \emph{avagāhate}.) A confluence? (\textbf{\sk{संगमः}}, \emph{saṅgamaḥ}.) Purifying? (\textbf{\sk{पावनः}}, \emph{pāvanaḥ}.) Sprinkles over? (\textbf{\sk{प्रोक्षति}}, \emph{prokṣati}.) A deep pool? (\textbf{\sk{ह्रदः}}, \emph{hradaḥ}.) Gives a bath to? (\textbf{\sk{स्नापयति}}, \emph{snāpayati}.) Divine? (\textbf{\sk{दिव्यः}}, \emph{divyaḥ}.) An arched gateway? (\textbf{\sk{तोरणम्}}, \emph{toraṇam}.) A pennant? (\textbf{\sk{पताका}}, \emph{patākā}.) Walks around? (\textbf{\sk{परिक्रामति}}, \emph{parikrāmati}.) Prayer beads? (\textbf{\sk{जपमाला}}, \emph{japamālā}.) A rudraksha bead? (\textbf{\sk{रुद्राक्षः}}, \emph{rudrākṣaḥ}.) A chariot procession? (\textbf{\sk{रथयात्रा}}, \emph{rathayātrā}.) Fetches? (\textbf{\sk{आहरति}}, \emph{āharati}.) A mendicant monk? (\textbf{\sk{भिक्षुः}}, \emph{bhikṣuḥ}.) A Buddhist monastery? (\textbf{\sk{विहारः}}, \emph{vihāraḥ}.) Alms? (\textbf{\sk{भिक्षा}}, \emph{bhikṣā}.) A lay follower? (\textbf{\sk{उपासकः}}, \emph{upāsakaḥ}.) A memorial shrine? (\textbf{\sk{चैत्यम्}}, \emph{caityam}.)
\end{tcolorbox}
